\section{Описание практической части}
\label{sec:Chapter4} \index{Chapter4}

В данной главе описывается программная реализация предложенного в главе~\ref{sec:Chapter3} решения: выбор языка и библиотек с обоснованием, организация кодовой базы, схема функционирования пайплайна, этап холодного старта (SFT), конфигурация обучения с подкреплением, теоретическая сложность ключевых алгоритмов и характеристики функционирования (скорость и память). Приводятся также фактические значения гиперпараметров, использованные в основной серии экспериментов.

\subsection{Язык реализации и выбор библиотек}
\label{subsec:impl-language}

Весь программный комплекс реализован на языке \textbf{Python}~3.11. Выбор обусловлен доминированием Python в экосистеме машинного обучения и наличием зрелых библиотек для всех подзадач работы: глубокого обучения, инференса LLM, работы с графами и обращения к внешним моделям через API. Зависимости фиксируются менеджером \texttt{poetry} (для \texttt{smart-thinking-llm}) и \texttt{setuptools}/\texttt{uv} (для \texttt{TunePPO}), что обеспечивает воспроизводимость окружения.

Ключевые библиотеки и мотивы их выбора:
\begin{itemize}
    \setlength{\itemsep}{\smallskipamount}
    \item \textbf{PyTorch} (\texttt{torch}, \texttt{torchao}) и \textbf{torchtune}~\cite{torchtune2024}~--- базовый фреймворк глубокого обучения и тонкой настройки LLM; служит вычислительной основой собственного RL-стека \texttt{TunePPO}.
    \item \textbf{VERL}~\cite{sheng2024verl}~--- промышленная платформа RL-обучения LLM, обеспечивающая эффективную интеграцию инференс-движка в цикл обучения и масштабируемое управление политикой, опорной моделью и наградой. Использована в качестве основной платформы для экспериментов на RuleTaker.
    \item \textbf{vLLM}~--- высокопроизводительный движок инференса с paged-attention и непрерывным батчингом; применяется для быстрого сэмплирования $G$ кандидатов на этапе rollout.
    \item \textbf{transformers}, \textbf{datasets}, \textbf{huggingface-hub}~--- загрузка предобученных моделей семейства Qwen3, токенизаторов с chat-шаблонами и работа с датасетами.
    \item \textbf{PEFT/LoRA}~\cite{hu2022lora} и \textbf{bitsandbytes}~--- параметр-эффективное дообучение (LoRA-адаптеры) и пейджинг состояния оптимизатора (\texttt{PagedAdamW}), снижающие потребление видеопамяти.
    \item \textbf{NetworkX}~--- представление графов рассуждений и вычисление Graph Edit Distance (раздел~\ref{subsec:metrics-formalization}).
    \item \textbf{Pydantic}~--- описание схем Structured Output для надёжного извлечения графов вспомогательной LLM (\texttt{StructuredGraphCreator}).
    \item \textbf{openai}~--- унифицированный клиент к OpenAI-совместимому API инференс-сервера (для генерации датасета и LLM-ассесора).
    \item \textbf{rapidfuzz}~--- быстрое вычисление нормализованного расстояния Левенштейна при сопоставлении сущностей в награде для Wikidata.
    \item \textbf{TensorBoard} и \textbf{Weights\,\&\,Biases}~--- логирование метрик качества и компонент награды в ходе оценки и обучения.
    \item \textbf{evalscope}~--- унифицированный запуск оценки моделей на внешних бенчмарках и на собственном адаптере RuleTaker.
\end{itemize}

\subsection{Организация кодовой базы}
\label{subsec:impl-structure}

Решение разделено на два репозитория, что отражает разделение ответственности, описанное в разделе~\ref{subsec:framework-architecture}:
\begin{itemize}
    \setlength{\itemsep}{\smallskipamount}
    \item \texttt{smart-thinking-llm} (около $4{,}7$~тыс. строк Python)~--- всё, что относится к данным, графам, метрикам и оценке: модули \texttt{datasets/}, \texttt{find\_patterns/}, \texttt{generation/}, \texttt{tools/graph\_creator/}, \texttt{tools/metrics/}, \texttt{prompts/}, а также инфраструктура оценки \texttt{evalscope/}.
    \item \texttt{TunePPO}~\cite{tuneppo2024}~--- RL-цикл: реализация графовых функций награды (\texttt{verl\_exp/\allowbreak edge\_rea\-soning\_reward.py} для RuleTaker и \texttt{GraphMultihop\allowbreak QAReward} в \texttt{ppotune/\allowbreak reward.py} для Wikidata), скрипты подготовки данных (\texttt{ruletaker\_\allowbreak prepare\_\allowbreak data\_kg.py}, \texttt{ruletaker\_\allowbreak prepare\_\allowbreak data\_struct.py}) и запуска обучения поверх VERL и torchtune.
\end{itemize}

Модульная организация (абстрактные базовые классы \texttt{BaseDataset}, \texttt{BasePatternFinder}, \texttt{BaseQuestionGenerator}, базовый \texttt{GraphCreator} и интерфейс \texttt{IRewardModel}) обеспечивает заменяемость компонентов: добавление нового датасета, экстрактора графа или функции награды не затрагивает остальные подсистемы, что выполняет нефункциональное требование модульности из раздела~\ref{subsec:requirements}. Общая структура комплекса и взаимодействие двух репозиториев показаны на рис.~\ref{fig:architecture}.

\begin{figure}[H]
\centering
\resizebox{\textwidth}{!}{%
\begin{tikzpicture}[
    mod/.style={draw, rounded corners, align=center, font=\footnotesize, minimum height=7mm, inner sep=3pt, fill=black!3},
    grp/.style={draw, dashed, rounded corners, inner sep=5mm},
    arr/.style={-{Stealth[length=2mm]}, thick},
    node distance=4mm,
]
% smart-thinking-llm modules
\node[mod] (data) {datasets/};
\node[mod, below=of data] (find) {find\_patterns/};
\node[mod, below=of find] (gen) {generation/};
\node[mod, right=14mm of data] (gc) {graph\_creator/};
\node[mod, below=of gc] (metrics) {metrics/};
\node[mod, below=of metrics] (eval) {evalscope/};
\begin{scope}[on background layer]
\node[grp, fit=(data)(find)(gen)(gc)(metrics)(eval), label={[font=\small\bfseries]above:smart-thinking-llm}] (stl) {};
\end{scope}
% TunePPO modules
\node[mod, right=30mm of gc] (reward) {награды:\\edge / multihop};
\node[mod, below=of reward] (prep) {prepare\_data};
\node[mod, below=of prep] (train) {VERL / torchtune};
\begin{scope}[on background layer]
\node[grp, fit=(reward)(prep)(train), label={[font=\small\bfseries]above:TunePPO}] (tppo) {};
\end{scope}
\draw[arr] (stl.east|-reward) -- node[above, font=\scriptsize]{$\hat{G}, G^*$} (reward.west);
\draw[arr] (gc) -- (metrics);
\draw[arr] (reward) -- (train);
\draw[arr] (prep) -- (train);
\end{tikzpicture}%
}
\caption{Архитектура программного комплекса: smart-thinking-llm (данные, графы, метрики, оценка) и TunePPO (награды и RL-обучение)}
\label{fig:architecture}
\end{figure}

\subsection{Схема функционирования}
\label{subsec:impl-flow}

Сквозной процесс работы комплекса включает следующие фазы.

\begin{enumerate}
    \setlength{\itemsep}{\smallskipamount}
    \item \textbf{Подготовка данных.} Скрипт \texttt{ruletaker\_prepare\_data\_kg.py} читает исходный RuleTaker, выполняет схлопывание вершин-правил, фильтрацию по отрицаниям и NatLang-подмножеству, формирует для каждого примера эталонное отображение \texttt{proof\_rules} (заключение $\mapsto$ списки допустимых наборов посылок), разбивает данные в соотношении $0{,}9:0{,}1$ с \texttt{seed}~$=42$ и сериализует в JSON-файлы \texttt{train.json}/\texttt{val.json} с полями \texttt{prompt}, \texttt{proof\_rules}, \texttt{label}.
    \item \textbf{Холодный старт (SFT).} Базовая модель адаптируется к формату \texttt{<think>...</think>\allowbreak<answer>...</answer>} (раздел~\ref{subsec:impl-sft}).
    \item \textbf{Rollout.} На каждом шаге обучения политика генерирует через vLLM по $G=8$ ответов на запрос с температурой $T=0{,}7$.
    \item \textbf{Извлечение графа и вычисление награды.} Для каждого ответа функция \texttt{compute\_score} разбирает теги, извлекает JSON-список рёбер, вычисляет покрытие эталонного подграфа процедурой \texttt{validate\_subgraph} и суммирует шесть компонент награды~(\ref{eq:reward-total}). Вычисление награды распараллелено (\texttt{reward.num\_workers}~$=16$).
    \item \textbf{Обновление политики.} GRPO нормализует преимущества внутри группы и обновляет LoRA-адаптеры.
    \item \textbf{Валидация и логирование.} Каждые $10$ шагов на отложенной выборке (до $1000$ примеров) рассчитываются метрики и покомпонентная статистика награды, логируемые в W\&B/TensorBoard; чекпойнты сохраняются каждые $10$ шагов.
\end{enumerate}

\subsection{Этап холодного старта (SFT)}
\label{subsec:impl-sft}

Перед RL-этапом реализован этап \emph{холодного старта} (cold-start Supervised Fine-Tuning), назначение которого~--- приучить компактную модель к жёсткому формату вывода и повысить устойчивость последующего RL-сигнала. Обучение запускается тренером \texttt{verl.trainer.sft\_trainer} на одной GPU (\texttt{torchrun --nproc\_per\_node=1}) поверх базовой модели \texttt{Qwen/Qwen3-0.6B} с LoRA-адаптерами ранга $r=64$, $\alpha=16$; целевой сигнал~--- максимизация правдоподобия корректного структурированного ответа (метка \texttt{entailment}/\texttt{not entailment}) на подготовленном датасете \texttt{KGReasoning}, одна эпоха обучения. Полученная модель (обозначаемая \texttt{Qwen3-0.6B-SFT-LoRA}) оценивается наравне с остальными конфигурациями и, кроме того, служит начальной политикой для комбинированного режима SFT\,$\to$\,RL (раздел~\ref{subsec:impl-sft-rl}).

Следует подчеркнуть методологическое разделение: \emph{контрастная} серия RL-экспериментов ($\alpha = 0$ против $\alpha > 0$), на которой основан главный вывод (раздел~\ref{subsec:impl-results}), инициализируется не SFT-чекпойнтом, а \emph{базовой} моделью Qwen3-0.6B. Это исключает влияние холодного старта на сравнение и позволяет приписать наблюдаемую разницу именно графовой компоненте награды. Дополнительно режим SFT\,$\to$\,RL, инициализируемый из SFT-чекпойнта, показывает наилучший результат и демонстрирует взаимодополняемость двух этапов.

\subsection{Конфигурация и инфраструктура RL-обучения}
\label{subsec:impl-rl-config}

Основная серия экспериментов на RuleTaker проведена на платформе VERL с алгоритмом GRPO~\cite{shao2024deepseekmath} в следующей конфигурации (одна GPU класса H100):
\begin{itemize}
    \setlength{\itemsep}{\smallskipamount}
    \item модель \texttt{Qwen/Qwen3-0.6B}, LoRA $r=64$, $\alpha=16$ на матрицах внимания и MLP; \texttt{flash\_attention\_2}, gradient checkpointing, тип данных \texttt{bfloat16};
    \item размер обучающего батча запросов $128$, число ответов на запрос $G=8$; максимальная длина запроса $1024$, ответа~--- $4096$ токенов; \texttt{max\_model\_len}~$=8192$;
    \item оптимизация: learning rate $1\cdot 10^{-4}$, \texttt{ppo\_mini\_batch\_size}~$=64$, \texttt{ppo\_micro\_batch\_size\_per\_gpu}~$=8$, динамический батчинг по токенам (\texttt{ppo\_max\_token\_len\_per\_gpu}~$=24576$);
    \item регуляризация: KL не вносится в награду (\texttt{use\_kl\_in\_reward=False}), вместо этого используется KL-член в функции потерь типа \texttt{low\_var\_kl} с коэффициентом $0{,}001$ и контроллер $\texttt{kl\_coef}=0{,}1$;
    \item rollout через vLLM: \texttt{tensor\_model\_parallel\_size}~$=1$, \texttt{gpu\_memory\_utili\-zation}~$=0{,}7$, \texttt{max\_num\_seqs}~$=256$, температура $0{,}7$;
    \item одна эпоха обучения, валидация и сохранение чекпойнтов каждые $10$ шагов.
\end{itemize}

Контрастность baseline/предложенного метода реализована единственным параметром награды \texttt{graph\_validity\_scale} ($=\alpha$): значение $0$ отключает графовую компоненту (baseline, только outcome-награда), значение $1$ включает плотную графовую награду. Прочие коэффициенты награды фиксированы (\texttt{correct\_answer\_reward}~$=100$, \texttt{graph\_validity\_reward}~$=100$, \texttt{graph\_parse\_bonus}~$=5$, \texttt{format\_penalty}~$=10$), что соответствует формуле~(\ref{eq:reward-total}).

Для экспериментов на Wikidata использован собственный стек \texttt{TunePPO} поверх torchtune с оптимизатором \texttt{PagedAdamW} (\texttt{bitsandbytes}); награда \texttt{GraphMultihop\allowbreak QAReward} приводит извлечённые сущности к каноническому виду вспомогательным экстрактором (Qwen3-32B) и сопоставляет их по нормализованному расстоянию Левенштейна с порогом $0{,}8$ в режиме \texttt{entity\_only}.

\subsection{Теоретическая сложность ключевых алгоритмов}
\label{subsec:impl-complexity}

\textbf{Граф\-метрики на множествах.} Метрики Node/Edge Precision, Recall и $F_1$ сводятся к операциям над множествами вершин и рёбер и вычисляются за $O(|\hat{V}| + |V^*| + |\hat{E}| + |E^*|)$.

\textbf{Graph Edit Distance.} Точное вычисление GED~--- NP-трудная задача; используемая реализация \texttt{networkx.graph\_edit\_distance} в худшем случае экспоненциальна по числу вершин. Однако в условиях задачи графы рассуждений малы (единицы~--- десятки вершин), поэтому GED вычисляется за приемлемое время; кроме того, для обучающего сигнала GED не применяется~--- там используется существенно более дешёвая процедура покрытия.

\textbf{Рекурсивная валидация подграфа.} Процедура \texttt{validate\_subgraph}, на которой основан обучающий сигнал, выполняет обход узлов-заключений начиная от вопроса с мемоизацией (\texttt{visited}), сравнивая для каждого заключения множество предсказанных посылок с эталонными вариантами. Её сложность линейна по размеру графа~--- $O(|\hat{E}| + |\hat{V}|)$, что критично, поскольку награда вычисляется для каждого из $G$ кандидатов на каждом шаге обучения.

\textbf{Сопоставление сущностей (Wikidata).} Канонизация и сопоставление по расстоянию Левенштейна выполняются за $O(n\cdot m\cdot L)$, где $n$~--- число предсказанных сущностей, $m$~--- число кандидатов, $L$~--- средняя длина строки; на практике доминирует обращение к вспомогательному экстрактору.

\subsection{Характеристики функционирования}
\label{subsec:impl-performance}

\textbf{Память.} Применение LoRA замораживает базовые веса и сокращает число обучаемых параметров примерно до $1\%$ от общего числа, что вместе с gradient checkpointing, вычислениями в \texttt{bfloat16} и пейджингом состояния оптимизатора позволяет дообучать модель Qwen3-0.6B на \emph{одной} GPU класса H100. Доля видеопамяти, резервируемая под vLLM-инференс, ограничена параметром \texttt{gpu\_memory\_utilization}~$=0{,}7$, оставляя место под градиенты и активации актора.

\textbf{Скорость.} Использование vLLM с непрерывным батчингом и динамической упаковки последовательностей по токенам обеспечивает высокую пропускную способность фазы rollout; параллельное вычисление награды ($16$ воркеров) исключает её из узких мест. Прямой JSON-разбор графа (подход~1 из раздела~\ref{subsec:graph-extraction}) на порядок быстрее, чем извлечение вспомогательной LLM, и потому выбран основным в RL-цикле, тогда как затратный \texttt{StructuredGraphCreator} включается лишь как fallback.

\textbf{Накладные расходы инициализации.} Построение индекса графа знаний Wikidata5m (загрузка троек и алиасов сущностей/отношений) занимает порядка $3$--$4$~минут при старте, после чего обращения к графу выполняются в оперативной памяти.

\subsection{Инфраструктура оценки}
\label{subsec:impl-eval}

Для воспроизводимой оценки реализован собственный адаптер RuleTaker в составе \texttt{evalscope} (\texttt{evalscope/benchmarks/ruletaker/ruletaker\_adapter.py}) с метрикой \texttt{rule\-taker\_accuracy}, а также конфигурация \texttt{evalscope\_multi\_runner} для пакетного прогона нескольких моделей. Параллельно модели оцениваются на внешних бенчмарках рассуждений (\texttt{math\_500}, \texttt{aime24}/\texttt{aime25}, \texttt{gpqa}, \texttt{mmlu\_redux}, \texttt{zebra\_logic}, \texttt{live\_code\_bench}), что позволяет контролировать сохранность общих способностей модели после дообучения на узкой логической задаче.

\textbf{Протокол оценки.} Оценка проводится в режиме \texttt{pass@1} (один сэмпл на запрос) с температурой $0{,}7$; на целевой задаче RuleTaker используется полный валидационный срез ($10\,000$ примеров). Для каждой модели вычисляются: точность предсказания метки \texttt{entailment}/\texttt{not entailment} (метрика \texttt{ruletaker\_accuracy}) и, по мере готовности соответствующего режима бенчмарка, структурные метрики (покрытие эталонного подграфа, Edge~$F_1$, GES). Модели, не приученные к строгому формату, оцениваются с извлечением графа подходом~2 (\texttt{StructuredGraphCreator}), модели семейства Qwen3~--- прямым разбором блока \texttt{<answer>}. Внешние бенчмарки запускаются с конфигурациями по умолчанию из \texttt{evalscope}, число прогонов варьируется, что объясняет разброс значений в таблице~\ref{tab:generalization}. Результаты приводятся и анализируются в разделе~\ref{subsec:impl-results}.

\subsection{Результаты экспериментов}
\label{subsec:impl-results}

Оценка проведена для пяти конфигураций модели Qwen3-0.6B на полном валидационном срезе RuleTaker ($10\,000$ примеров): базовая модель, модель после холодного старта (SFT), две контрастные RL-политики, обученные из базовой модели с бинарной ($\alpha=0$) и графовой ($\alpha>0$) наградой, а также комбинированный двухэтапный режим SFT\,$\to$\,RL с графовой наградой.

\textbf{Точность на целевой задаче (RuleTaker).} Точность предсказания метки \texttt{entailment}/\texttt{not entailment} приведена в таблице~\ref{tab:ruletaker-acc}.

\begin{table}[htbp]
\centering
\caption{Точность на RuleTaker ($N=10\,000$)}
\label{tab:ruletaker-acc}
\begin{tabular}{|l|c|c|}
\hline
\textbf{Конфигурация} & \textbf{Награда} & \textbf{Accuracy} \\
\hline
Qwen3-0.6B (база)        & ---                    & $0{,}6715$ \\
\hline
Qwen3-0.6B-SFT-LoRA      & SFT (max-likelihood)   & $0{,}6923$ \\
\hline
Qwen3-0.6B-RL-baseline   & бинарная, $\alpha=0$   & $0{,}6991$ \\
\hline
Qwen3-0.6B-RL            & графовая, $\alpha>0$   & $0{,}7115$ \\
\hline
Qwen3-0.6B-SFT-RL        & SFT\,+\,графовая, $\alpha>0$ & $\mathbf{0{,}7143}$ \\
\hline
\end{tabular}
\end{table}

Точность растёт монотонно по мере усложнения обучающего сигнала: $0{,}6715 \to 0{,}6923 \to 0{,}6991 \to 0{,}7115 \to 0{,}7143$. Существенно, что графовая награда ($\alpha>0$) превосходит бинарный baseline ($\alpha=0$) на $1{,}2$ процентных пункта ($0{,}7115$ против $0{,}6991$), а наилучший результат показывает двухэтапный режим SFT\,$\to$\,RL. Таким образом, плотная графовая награда повышает точность даже на бинарной целевой задаче.

\textbf{Структурные метрики.} Главный эффект графовой награды проявляется в структурном качестве цепочек рассуждений (таблица~\ref{tab:struct-metrics}), измеренном на режиме \texttt{ruletaker\_graph} бенчмарка \texttt{ruletaker\_adapter}.

\begin{table}[htbp]
\centering
\caption{Структурные метрики цепочек рассуждений на RuleTaker ($N=10\,000$)}
\label{tab:struct-metrics}
\begin{tabular}{|l|c|c|c|}
\hline
\textbf{Конфигурация} & \textbf{Node $F_1$} & \textbf{Edge $F_1$} & \textbf{GES} \\
\hline
Qwen3-0.6B (база)        & $0{,}2117$ & $0{,}1901$ & $0{,}8354$ \\
\hline
Qwen3-0.6B-SFT-LoRA      & $0{,}2362$ & $0{,}2174$ & $0{,}8453$ \\
\hline
Qwen3-0.6B-RL-baseline ($\alpha=0$) & $0{,}2251$ & $0{,}1875$ & $0{,}8475$ \\
\hline
Qwen3-0.6B-RL ($\alpha>0$) & $0{,}2591$ & $0{,}2335$ & $0{,}8525$ \\
\hline
Qwen3-0.6B-SFT-RL ($\alpha>0$) & $\mathbf{0{,}2671}$ & $\mathbf{0{,}2394}$ & $\mathbf{0{,}8563}$ \\
\hline
\end{tabular}
\end{table}

Здесь различие между подходами выражено ярко. Бинарный baseline, оптимизирующий только финальный ответ, \emph{не улучшает} структуру вывода: его Edge~$F_1$ ($0{,}1875$) даже немного ниже, чем у базовой модели ($0{,}1901$) и у SFT ($0{,}2174$),~--- повышая точность метки, outcome-награда не побуждает модель строить корректную цепочку. Напротив, графовая награда повышает Node~$F_1$ до $0{,}2591$ и Edge~$F_1$ до $0{,}2335$, а двухэтапный режим SFT\,$\to$\,RL~--- до $0{,}2671$ и $0{,}2394$ соответственно; GES растёт монотонно ($0{,}8354 \to 0{,}8563$). На выборке из $10\,000$ примеров эти различия устойчивы и согласованы по всем трём структурным метрикам, что подтверждает центральную гипотезу работы: графовая награда улучшает структурную корректность рассуждений, недостижимую при обучении только по результату.

\textbf{Сохранность общих способностей.} Поскольку дообучение проводилось на узкой логической задаче, существует риск деградации общих рассуждающих способностей модели (catastrophic forgetting). Для контроля этого эффекта базовая модель, модель после холодного старта (\texttt{Qwen3-0.6B-SFT-LoRA}) и обе RL-политики оценены на внешних бенчмарках; сводка приведена в таблице~\ref{tab:generalization}.

\begin{table}[htbp]
\centering
\caption{Оценка на внешних бенчмарках (mean accuracy; для \texttt{zebra}~--- puzzle accuracy)}
\label{tab:generalization}
\small
\begin{tabular}{|l|c|c|c|c|}
\hline
\textbf{Конфигурация} & \textbf{math\_500} & \textbf{gpqa} & \textbf{mmlu\_redux} & \textbf{zebra} \\
\hline
Qwen3-0.6B (база)        & $0{,}738$ & $0{,}227$ & $0{,}571$ & $0{,}267$ \\
\hline
Qwen3-0.6B-SFT-LoRA      & $0{,}744$ & $0{,}268$ & $0{,}561$ & $0{,}263$ \\
\hline
Qwen3-0.6B-RL-baseline   & $0{,}760$ & $\mathbf{0{,}293}$ & $0{,}562$ & $0{,}273$ \\
\hline
Qwen3-0.6B-RL            & $0{,}746$ & $0{,}263$ & $0{,}566$ & $0{,}286$ \\
\hline
Qwen3-0.6B-SFT-RL        & $\mathbf{0{,}766}$ & $0{,}278$ & $0{,}565$ & $\mathbf{0{,}295}$ \\
\hline
\end{tabular}
\end{table}

Дообучение на узкой логической задаче RuleTaker не привело к деградации общих способностей: на \texttt{math\_500}, \texttt{gpqa\_diamond}, \texttt{mmlu\_redux} и \texttt{zebra\_logic} все дообученные конфигурации остаются на уровне базовой Qwen3-0.6B или превосходят её. Более того, на логической головоломке \texttt{zebra\_logic} наблюдается последовательный рост ($0{,}267 \to 0{,}286 \to 0{,}295$ для базы, RL и SFT\,$\to$\,RL), что указывает на положительный перенос навыка структурного рассуждения за пределы обучающего датасета. Таким образом, предложенная графовая награда обучает структуре рассуждений, не разрушая~--- а в ряде случаев усиливая~--- общую способность модели к выводу.

\subsection{Двухэтапный пайплайн SFT\,+\,RL}
\label{subsec:impl-sft-rl}

Полная методология предусматривает два последовательных этапа обучения, схематически показанных на рис.~\ref{fig:sft-rl}: холодный старт (SFT) адаптирует базовую модель к формату структурированных рассуждений, после чего этап RL с графовой наградой оптимизирует структурную корректность цепочек вывода.

\begin{figure}[H]
\centering
\resizebox{\textwidth}{!}{%
\begin{tikzpicture}[
    stage/.style={draw, rounded corners, align=center, minimum height=12mm, minimum width=34mm, font=\small, fill=black!4},
    mod/.style={draw, align=center, minimum height=12mm, minimum width=26mm, font=\small, fill=black!8},
    arr/.style={-{Stealth[length=2mm]}, thick},
    node distance=10mm,
]
\node[mod] (base) {Qwen3-0.6B\\\footnotesize (base)};
\node[stage, right=of base] (sft) {Этап 1: SFT\\\footnotesize cold-start, max-likelihood};
\node[mod, right=of sft] (ckpt) {Qwen3-0.6B\\-SFT-LoRA};
\node[stage, right=of ckpt] (rl) {Этап 2: RL (GRPO)\\\footnotesize графовая награда};
\node[mod, right=of rl] (final) {Обученная\\политика};
\draw[arr] (base) -- (sft);
\draw[arr] (sft) -- (ckpt);
\draw[arr] (ckpt) -- (rl);
\draw[arr] (rl) -- (final);
\draw[arr, dashed] (base.south) to[out=-30, in=-150] node[below, font=\scriptsize, align=center]{контрастная серия: RL напрямую от базы ($\alpha{=}0$ vs $\alpha{>}0$)} (rl.south);
\end{tikzpicture}%
}
\caption{Двухэтапный пайплайн обучения. Пунктиром показана контрастная серия, инициализируемая базовой моделью в обход SFT}
\label{fig:sft-rl}
\end{figure}

Сравнение всех пяти конфигураций (таблицы~\ref{tab:ruletaker-acc} и~\ref{tab:struct-metrics}) показывает, что вклады двух этапов складываются. Холодный старт (SFT) уже повышает как точность ($0{,}6715 \to 0{,}6923$), так и структурные метрики (Edge~$F_1$: $0{,}1901 \to 0{,}2174$), приучая модель к корректному формату цепочки. Этап RL с графовой наградой, инициализированный из базовой модели, даёт дополнительный прирост по структуре (Edge~$F_1$ $0{,}2335$, Node~$F_1$ $0{,}2591$). Наконец, комбинированный режим SFT\,$\to$\,RL достигает наилучших значений по всем метрикам одновременно (Accuracy $0{,}7143$, Node~$F_1$ $0{,}2671$, Edge~$F_1$ $0{,}2394$, GES $0{,}8563$), подтверждая, что холодный старт и графовая награда дополняют друг друга, а не конкурируют.

Показательно сравнение двух RL-конфигураций, обученных из \emph{одной} отправной точки: бинарный baseline ($\alpha=0$) повышает точность, но оставляет структурное качество на уровне базовой модели, тогда как графовая награда ($\alpha>0$) улучшает и то, и другое. Это изолирует вклад именно графовой компоненты награды.

\subsection{Ограничения реализации и угрозы валидности}
\label{subsec:impl-limitations}

Для корректной интерпретации результатов отметим ограничения реализации и потенциальные угрозы валидности выводов.
\begin{itemize}
    \setlength{\itemsep}{\smallskipamount}
    \item \textit{Отсутствие проверки статистической значимости.} Сравнение конфигураций (таблицы~\ref{tab:ruletaker-acc}, \ref{tab:struct-metrics}) проведено на большой выборке ($10\,000$ примеров), и различия согласованы по всем структурным метрикам, однако формальная проверка статистической значимости и оценка устойчивости по нескольким случайным инициализациям выходят за рамки работы.
    \item \textit{Зависимость от качества извлечения графа.} Точность графовой награды ограничена качеством парсинга ответа: ошибки извлечения троек (особенно у компактной модели на ранних шагах обучения) вносят шум в обучающий сигнал; этот эффект частично компенсируется штрафами за нарушение формата и fallback-экстрактором.
    \item \textit{Масштаб модели и единственная архитектура.} Эксперименты проведены на одной компактной модели (Qwen3-0.6B); обобщение выводов на другие размеры и семейства моделей требует отдельной проверки.
    \item \textit{Разброс на внешних бенчмарках.} Значения в таблице~\ref{tab:generalization} получены при разном числе прогонов и малых выборках (особенно \texttt{aime}), поэтому их следует интерпретировать как индикативные, а не как точные оценки.
    \item \textit{Специфика сопоставления для Wikidata.} Сопоставление сущностей по нормализованному расстоянию Левенштейна с фиксированным порогом $0{,}8$ может приводить к ложным совпадениям/пропускам для близких по написанию сущностей.
\end{itemize}
Учёт этих ограничений определяет направления дальнейшей работы (раздел~\ref{subsec:conclusion-future}).

\subsection{Воспроизводимость и программное окружение}
\label{subsec:impl-repro}

Для обеспечения воспроизводимости все зависимости зафиксированы: окружение \texttt{smart-thinking-llm} управляется \texttt{poetry} (с экспортом в \texttt{requirements.txt}), окружение \texttt{TunePPO}~--- спецификацией \texttt{pyproject.toml}. Версии ключевых библиотек приведены в таблице~\ref{tab:versions}. Случайность контролируется фиксированным значением \texttt{seed}~$=42$ при разбиении данных и сэмплировании. Для управления крупными артефактами данных (граф Wikidata5m, подготовленные датасеты) используется \texttt{DVC}. Запуск обучения оформлен в виде shell-скриптов (\texttt{h100\_edge\_run.sh}, \texttt{struct\_run.sh}), фиксирующих все гиперпараметры командной строки.

\begin{table}[H]
\centering
\caption{Версии ключевых библиотек программного окружения}
\label{tab:versions}
\small
\begin{tabular}{|l|c|l|}
\hline
\textbf{Библиотека} & \textbf{Версия} & \textbf{Назначение} \\
\hline
Python        & 3.11        & язык реализации \\
torchtune     & 0.6.1       & RL-стек TunePPO \\
vLLM          & 0.8.5       & инференс на rollout \\
transformers  & 4.49+       & модели и токенизаторы \\
networkx      & 3.5         & графы и GED \\
pydantic      & 2.11        & схемы Structured Output \\
openai        & 1.97        & клиент инференс-API \\
rapidfuzz     & 3.12        & расстояние Левенштейна \\
datasets      & 3.2         & загрузка датасетов \\
wandb         & 0.21        & логирование обучения \\
tensorboard   & 2.20        & логирование оценки \\
\hline
\end{tabular}
\end{table}

Таким образом, сочетание фиксированных версий, скриптов запуска с явными гиперпараметрами, версионирования данных через DVC и публикации датасета \texttt{KGReasoning} обеспечивает полную воспроизводимость экспериментов, что соответствует требованию документированности из раздела~\ref{subsec:requirements}.
